% =====================================================================
%  Pacotes e configuracoes gerais do preambulo.
% =====================================================================

% --- Codificacao e idioma --------------------------------------------
\usepackage[T1]{fontenc}
\usepackage[utf8]{inputenc}
\usepackage[brazil]{babel}
\usepackage{lmodern}

% --- Elementos graficos ----------------------------------------------
\usepackage{graphicx}
\graphicspath{{figures/}}
\usepackage{booktabs}
\usepackage{multirow}

% --- Matematica ------------------------------------------------------
\usepackage{amsmath}
\usepackage{amssymb}

% --- Tipografia ------------------------------------------------------
\usepackage{microtype}      % ajuste fino de espacamento e protrusao
\usepackage{xcolor}
\usepackage{epigraph}       % epigrafes de capitulo

% --- Indice remissivo ------------------------------------------------
% Marque termos no texto com \index{termo}; o indice se monta sozinho.
\makeindex

% --- Links -----------------------------------------------------------
\usepackage{hyperref}
\definecolor{lkLink}{RGB}{0,70,140}
\hypersetup{
  colorlinks=true,
  linkcolor=lkLink,
  citecolor=lkLink,
  urlcolor=lkLink,
  bookmarksdepth=3
}
% Para a versao de impressao, troque a linha acima por:
%   \hypersetup{hidelinks}

% --- Comandos de conveniencia ----------------------------------------
% Fonte de uma figura ou tabela.
\newcommand{\fonte}[1]{%
  \par\vspace{0.2em}{\footnotesize Fonte: #1}%
}
